
\section{Related Work}

\textbf{Prototype-based meta-learning.} 
Prototype-based meta-learning is based on distance metrics and generally learns a shared or adaptive embedding space in which query images are accurately matched to support images for classification. It relies on the assumption that a common metric space is shared across related tasks and usually does not employ an explicit base learner for each task. By extending the matching network~\cite{vinyals2016matching} to few-shot scenarios, Snell \etal~\cite{snell2017prototypical} constructed a prototype for each class by averaging the feature representations of samples from the class in the metric space. The classification matches the query samples to prototypes by computing their distances. To enhance the prototype representation, Allen \ et al~\cite {allen2019infinite} proposed an infinite mixture of prototypes to adaptively represent data distributions for each class, using multiple clusters instead of a single vector. Oreshkin \ et al~\cite {oreshkin2018tadam} proposed a task-dependent adaptive metric for few-shot learning and established prototype classes conditioned on a task representation encoded by a task embedding network. Yoon \ et al~\cite {yoon2019tapnet} proposed a few-shot learning algorithm aided by a linear transformer that performs task-specific null-space projection of the network output. Graphical neural network-based models generalize the matching methods by learning the message propagation from the support set and transferring it to the query set~\cite{garcia2018few}. Prototype-based methods have recently been improved in various ways \cite{cao2019theoretical,triantafillou2019meta,zhen2020memory}. To the best of our knowledge, we are the first to propose a diffusion method to generate the prototype per task, rather than a deterministic function.


\textbf{Diffusion models.} These models belong to a category of neural generative models that utilize stochastic diffusion processes \cite{sohl2015deep,song2020denoising}, much like those found in thermodynamics. In this framework, a gradual introduction of noise 
is applied to a sample of data. 
A neural model then learns to reverse the process by progressively removing noise from the sample. The model denoises an initial pure noise to obtain samples from the learned data distribution. Ho \etal~\cite{ho2020denoising} and Song \etal~\cite{song2020denoising} contributed to advancements in image generation, while Dhariwal and Nichol \cite{dhariwal2021diffusion} introduced classifier-guided diffusion for a conditioned generation. GLIDE later adapted this approach~\cite{nichol2021glide}, enabling conditioning on textual CLIP representations. Classifier-free guidance~\cite{ho2022classifier} allows for conditioning with a balance between fidelity and diversity, resulting in improved performance ~\cite{nichol2021glide}. Since the guided diffusion model requires a large number of image-annotation pairs for training, Hu \etal~\cite{hu2022self} propose self-guided diffusion models. 
Recently, Hyperdiffusion~\cite{lutati2022ocd, erkocc2023hyperdiffusion} was proposed in the weight space for  generating implicit neural fields and 3D reconstruction.
In this paper, we introduce ProtoDiff, a prototype-based meta-learning approach within a task-guided diffusion model that incrementally improves the prototype's expressiveness by utilizing a limited number of samples.


